%% ============================================================
\section{\sys{}: Research with Verifiability}
\label{sec:system}
%% ============================================================

We now describe \sys{}, an end-to-end autonomous research system whose three-stage architecture is shaped by the CoE requirements: each module is designed to produce structured artifacts that carry the provenance metadata needed to verify claims against their evidence (Figure~\ref{fig:architecture}).

\begin{figure}[t]
\centering
\includegraphics[width=\textwidth]{figures/00-overall-system.png}
\caption{\sys{} pipeline. Stage~1 grounds the literature via retrieved PDFs; Stage~2 explores and evaluates solutions across parallel branches; Stage~3 writes and verifies the paper, with a Claim Verifier that checks every claim against its evidence source before the final output is produced.}
\label{fig:architecture}
\end{figure}


\subsection{Stage 1: Literature Grounding}
\label{sec:system:pi}

The Problem Investigator (PI) is designed to ensure that every paper the system cites was retrieved from a scholarly database, read in full text, and recorded with provenance metadata.
Without structured retrieval, autonomous systems tend to generate citations from model memory---in our audit, systems without retrieval-grounded references exhibit hallucinated reference rates of up to 21\% (\S\ref{sec:exp:integrity}).
PI addresses this by construction: starting from seed papers, it builds a citation graph via scholarly database queries, reads up to 100 full-text PDFs per topic, and produces a structured research brief.
The brief feeds the Ideator, and PI's seed reference bibliography provides grounding material for citation claims in the final paper.
Pipeline details are in Appendix~\ref{app:system_details}.


\subsection{Stage 2: Discovery}
\label{sec:system:discovery}

The Ideator generates candidate approaches based on the PI brief, scores them on novelty and feasibility, and distributes the top-ranked proposals across parallel branches of the \emph{Parallel Explore-Exploit} (PEE) orchestrator.
Each branch runs an isolated cycle: a Solver agent iterates up to $E$ evaluated versions per node, with a task-specific evaluator scoring each submission.
At each iteration, the top-$K$ branches are retained, and the remaining slots are filled with new branches derived from these top performers via fresh ideation.
After $I$ iterations across $B$ branches, a best-run selector filters out solutions flagged for specification violations (Section~\ref{appx:i2_spec_violation}), selects the highest-scoring remaining solution, and runs ablation experiments on it.
The evaluator scores, execution logs, and ablation results are passed to Stage~3 as source material for paper writing and claim verification.
Architecture details are in Appendix~\ref{app:system_details}.


\subsection{Stage 3: Paper Writing \& Verification}
\label{sec:system:stage3}

\paragraph{Paper Writer.}
\label{sec:paper_writer}
\label{sec:system:writer}

The Paper Writer produces \LaTeX{} through a five-stage claim-grounded pipeline.
\textsc{Conceive} reads all assembled raw materials---PI brief, experimental log, verified scores, solver code, and seed-paper abstracts---and emits a \emph{research representation}: a markdown narrative where every factual claim carries an inline evidence tag binding it to a specific workspace artifact (a log line number, a score file entry, a citation key, or an ablation result).
\textsc{Ground} then validates each tag deterministically: the reported score must match the best-run score from discovery, baselines must be traceable to PI brief entries or marked \textsc{estimated}, and every referenced artifact must exist.
\textsc{Critic} audits what deterministic checks cannot---gap--approach alignment, internal contradictions, overclaims, missing comparisons, and baseline fairness---returning \textsc{pass} or a list of issues.
\textsc{Resolve} rewrites the representation against the \textsc{Ground} flags and \textsc{Critic} issues jointly, dropping unsupported claims and calibrating overclaims. The \textsc{Ground}--\textsc{Critic}--\textsc{Resolve} loop iterates until convergence or plateau.
Finally, \textsc{Compose} renders the grounded representation into \LaTeX{} one section at a time. Because each section writer receives verified numbers and named baselines alongside the representation, it writes prose around established facts rather than generating claims that must be sourced after the fact.

\paragraph{Claim Verifier and Refinement.}
\label{sec:claim_verifier}
\label{sec:system:verifier}

Even after grounding, the composed \LaTeX{} can introduce unsupported claims---through paraphrasing drift, misattributed citations, or numerical rounding errors.
The Claim Verifier catches these by checking every claim in the draft against its declared evidence source, dispatching on claim type: numerical claims against evaluator logs, citation claims against the bibliography with LLM-judged abstract entailment, and methodological claims against experimental logs.
Unsourced claims are flagged automatically.
A refinement pass then consumes the verifier's findings: an LLM rewrites flagged sentences to match their evidence sources, removes claims that cannot be supported, and strips all inline evidence annotations from the final \LaTeX{}.
Only a draft with no remaining blocking violations is promoted to the final paper.

